\section{A tower on the residual field}\label{b:scale-tower}

We apply the discrepancy calculus to the residual data
$S,H,A^+$ of \eqref{eq:A9}--\eqref{eq:A12}.  All varieties in this
section are over $\C$.  The construction is performed along the
sequence fixed there.  The symbols $\lesssim$ and $\asymp$ mean inequalities
with constants independent of its member.  Constants may depend on the
original bounded dimension and index data and on constants already fixed
at an earlier level.  There will be at most $\dim S+1$ levels; passing to
subsequences or discarding finitely many members at each level is therefore
harmless.  A positive rational described as fixed is independent of the
member, whereas extraction models and sufficiently small perturbations
used only to construct them may depend on that member.

At a collapsing level, we isolate an lc divisor carrying the smallest
scale and construct a contraction for which it is the only horizontal
component of the positive support of $H$. Weak essentiality bounds the
contraction threshold from below; bounded fibre degrees and a birational
section system bound it from above. The same degree comparison then
excludes every other positive horizontal component.

\subsection{The data carried by one level}

The discrepancy functions and the divisor $H$ define the contraction step.
The invariant-prime inequalities retain the regularity tests for sections.
We also keep a model of the initial residual field: its principal moduli
will bound the index after pushing to each successive base, without
accumulating fibre multiplicities.
We maintain the following data on a projective normal variety $S$, whose
function field is denoted by $K$.
\begin{enumerate}
\item \emph{The lc presentation.} $H$ is an effective ample rational Cartier divisor.  The discrepancy
function $A^+$ has an effective generalized lc CY presentation
$a_K+\mathbf M_0$, with boundary $B_0\ge H$ and a bounded b-Cartier
multiple of $\mathbf M_0$.
\item \emph{The positive discrepancy inherited from the preceding level.}
Except at the first level, there is an effective generalized
$\delta$-lc presentation $Q$, for a fixed $\delta>0$, with rational
b-Cartier b-nef part $\mathbf M_Q$, boundary $B_Q$, and total adjoint
$H/\lambda$, where
$\lambda>0$ is rational and tends to zero.  Initially $Q$ is absent and
$\lambda=0$.
\item \emph{The scale.} Write
\[
 \Lambda=H+\lambda Q,
 \qquad s=\min\bigl(\{1\}\cup
       \{\Lambda(P): A^+(P)=0\}\bigr).
\]
The primes in this definition are divisorial valuations anywhere over
$S$; the minimum is $1$ when there are none.  It exists and is positive:
for each member the rational presentations give a discrete group of
values; away from the first level $\Lambda\ge\delta\lambda$ on these
primes, while initially $A^++H=A_B>0$.  In particular $s\gtrsim\lambda$.
\item \emph{The section tests.} For every invariant prime $E$ of $\Spec R$ whose nontrivial
restriction to $K$ is $jP$, we have
\begin{equation}\label{eq:B2}
 jA^+(P)\le1-h_E,\qquad jH(P)\le h_E,\qquad
 j\Lambda(P)\le\lambda+C_0h_E.
\end{equation}
Every strict prime $P$ of $S$ with $H(P)=0$ has such a witness $E$ with
$h_E=0$ and bounded $j$.  A positive strict coefficient $H(P)$ is at
least a fixed $c_0>0$ if $A^+(P)>0$, and at least $c_0s$ if $A^+(P)=0$.
\item \emph{Essentiality and the initial field.} Weak essentiality holds in the form
\begin{equation}\label{b:scale-essentiality}
 P(H')\le P(H)+\eta\Lambda(P)
 \quad\text{if }H'\ge0,\ H'\sim_{\Q}H,\ A^+(P)=0,
 \qquad \eta\longrightarrow0.
\end{equation}
There is also a projective $\Q$-factorial Fano type model
$\mathcal Y$ of the initial field $K_1$, with a contraction to $S$,
such that $A^+$ is the valuative push of $A_1^+$ and $A_1^+$ has
effective trace on $\mathcal Y$.
\end{enumerate}
Here $A_1^+$ is the ordinary lc discrepancy with
$K_{S_1}+B_{0,1}=\frac1n\Div(f_h)$, equivalently with principal CY moduli of
denominator dividing $n$.  Initially the asserted conditions follow
from \eqref{eq:A9}--\eqref{eq:A12}.  For the strict coefficient bounds,
the witnesses there satisfy $jH(P)=h_E\in\frac1n\Z$; $j$ is bounded
when $A^+(P)>0$ by the discrepancy index gap.  For a strict lc prime,
$H(P)\ge s$.  Use a small $\Q$-factorialization for $\mathcal Y$;
the initial variety is of Fano type.

Pass to a subsequence on which either $s\to0$, called a collapsing
level, or $s$ stays bounded below.  In the latter case this is the last
level and we replace $s$ by $1$ in subsequent formulas, retaining the
actual positive lower bound when choosing constants and decreasing
the strict coefficient-gap constants accordingly.
\subsection{A positive perturbation and a birational system}

Choose fixed small positive rationals $a,b$, with $b$ sufficiently small
relative to the preceding constants, and put
\begin{equation}\label{b:scale-perturbation}
 \gamma=\frac{b\lambda}{s},\qquad
 A'=(1-\gamma)A^++\gamma Q+aH,\qquad d'=\frac bs-a.
\end{equation}
At the first level the convention is $\gamma Q=0$ and $d'=-a$.
The total adjoint of $A'$ is $d'H$, and its b-nef part is
$\mathbf M'=(1-\gamma)\mathbf M_0+\gamma\mathbf M_Q$.
Since $\lambda/s$ is bounded, $\gamma$ is uniformly small.  The trace is
effective because $B_0\ge H$ and both original boundaries are effective.
The datum is generalized klt: at later levels use $\gamma Q>0$, and at
the first use $A'=(1-a)A^++aA_B>0$.

\begin{lemma}\label{b:scale-bad-primes}
For a sufficiently small fixed $e>0$, every prime with $A'(P)<e$ is
lc for $A^+$.  There are finitely many such primes for each member.
At a collapsing level they satisfy
\[
                 \frac s2\le H(P)\le\frac ea.
\]
At the last level the constants may be chosen so there are none.
One can extract exactly their exceptional members on a projective
$\Q$-factorial Fano type model $S'\to S$.  At a collapsing level
one may additionally extract a prime $G_*$ with $A^+(G_*)=0$ and
$\Lambda(G_*)=s$.
On this model the trace of $A'$ is effective and $B_0\ge H$ still holds,
where $H$ now denotes its pullback.
\end{lemma}

\begin{proof}
The bounded b-index of $\mathbf M_0$ gives a uniform positive gap for
nonzero values of $A^+$.  Choose $e$ below this gap multiplied by
$1-\gamma$.  For a bad prime we therefore have $A^+(P)=0$, and
\[
 aH(P)<e,\qquad \lambda Q(P)<\frac{es}{b}
\]
when $Q$ is present.  Since $H(P)+\lambda Q(P)\ge s$, choosing
$e<b/2$ gives the asserted lower bound; initially it follows directly
from $H(P)\ge s$.  At the last level, if the unmodified scale is at
least $s_0>0$, these same formulas give a uniform positive lower bound
for $A'$ on lc primes, in terms of $a,b,s_0$.  Decrease $e$ further.

For completeness, the finiteness assertion uses more than the finiteness
of components on one resolution.  On a smooth b-Cartier descent of any
generalized klt datum, let $\Sigma$ be the reduced snc support of its
boundary.  Its discrepancy at $P$ is
\[
 A_{(W,\Sigma)}(P)+\sum_D a_D\ord_P(D),\qquad a_D>0.
\]
The first term is a nonnegative integer.  If the total is less than
$e<1$, this term is zero.  We recall why this forces $P$ to be monomial
at a stratum of $\Sigma$.  At the generic point of its centre choose
regular parameters $x_1,\ldots,x_\ell$, including the boundary normals,
and units $y_j$ lifting a separating transcendence basis of the centre.
The reduced-pair discrepancy is the discrepancy of
\[
 d\log x_1\wedge\cdots\wedge d\log x_\ell
       \wedge\bigwedge_jdy_j
\]
plus the orders of those $x_i$ which are not boundary normals.
On a smooth model carrying $P$, write $x_i=t^{w_i}u_i$, with $t$ a
uniformizer and $u_i$ units.  The displayed form has at most a simple
pole.  Equality zero thus forces every $x_i$ to be a boundary normal
and the logarithmic wedge to have nonzero residue.  The centre is a
stratum, and the wedge of the initials
$\overline u_i\overline t^{w_i}$ and $\overline y_j$ is nonzero.
In characteristic zero these initials are algebraically independent
in the required normal directions, also over the centre's function
field, which is algebraic over $\C(\overline y_j)$.  Expand an element
of the local ring to arbitrarily high finite order in the $x_i$, with
coefficients lifted from its residue field.  Independence prevents
cancellation in its leading polynomial, while the remainder has
arbitrarily large value.  Its value is therefore the minimum of its
monomial values.  The stratum and the positive integral weights $w_i$
determine $P$.  There are finitely many strata, and
$\sum a_Dw_D<1$ bounds all these weights.  This proves finiteness.

In the collapsing case choose $G_*$ attaining the scale, and write
$H_*=H(G_*)$.  Then
\[
                H_*\le s,\qquad A'(G_*)\le b+as.
\]
To extract the indicated finite set, replace $\gamma$ temporarily by
a much smaller positive rational $\widetilde\gamma<a\lambda$,
leaving it zero at the first level.  The resulting effective generalized
klt datum has ample anti-adjoint.  Its discrepancy at every requested
prime is less than $1$: these primes are lc for $A^+$, their
$H$-values have the bounds just proved, and
$\widetilde\gamma$ can be chosen for this finite set in the given
member.  Add a sufficiently small positive multiple of $H$ to both
moduli and total adjoint.  This leaves discrepancies unchanged, makes
the moduli big, and preserves ample anti-adjoint.  By
Lemma~\ref{b:ordinary-approximation}, replace the moduli by ordinary
effective klt data with no larger discrepancies.  The resulting
ordinary pair is log Fano.  Lemma~\ref{b:controlled-extraction}
extracts precisely the requested exceptional primes on a
$\Q$-factorial Fano type model.
For the original $A'$, their discrepancies are either less than $e$
or at most $b+as<1$, so its trace remains effective.  At these primes
$B_0(P)=1$ and $H(P)\le\max(e/a,s)\le1$ after the indicated choices.
Consequently $B_0\ge H$ also survives the extraction.
\end{proof}

\begin{lemma}\label{b:scale-birationality}
At every level there is a positive integer $k\le C_1/s$ such that
\begin{equation}\label{eq:B3}
                         |kH|\text{ is birational on }S.
\end{equation}
For a rational divisor this denotes the space of rational functions
$f$ with $\Div(f)+kH\ge0$; no bounded Cartier index for $H$ is asserted.
\end{lemma}

\begin{proof}
There is nothing to prove in dimension zero.  First extract just the
bad primes by Lemma~\ref{b:scale-bad-primes}, without an additional
pivot.  Every positive strict coefficient of $H$ is now at least a
fixed positive multiple of $s$.  Moreover the $A'$-boundary coefficient
at each such prime has a uniform positive lower bound.  At a new prime
it is greater than $1-e$.  At an old non-lc prime use
$B_0(P)\ge H(P)\ge c_0$ and the effectiveness of the other boundary;
at an old lc prime use $B_0(P)=1$ and $H(P)\le1$.

Let $G$ be the sum of the bad primes, which are now strict.  With a
sufficiently large fixed $C_2$ and $k_0=\lceil C_2/s\rceil$, the integral
Weil divisor
\[
                  N=\lfloor k_0H\rfloor-G
\]
satisfies $N\ge k_0H/2$.  Indeed the rounding error plus the coefficient
of $G$ is at most $2$ at each prime of the support, whereas each
positive coefficient of $k_0H$ is uniformly large.
Choose a fixed sufficiently small rational $\sigma>0$.  Subtract
$\sigma(\{k_0H\}+G)$ from the $A'$-boundary and add
$(\sigma k_0-d')\overline H$ to its moduli; choose $C_2$ large after
$\sigma$ so this last coefficient is positive.  The total adjoint is
$\sigma N$, the boundary is effective, and its discrepancy is
\[
             A'+\sigma\overline{(\{k_0H\}+G)}.
\]
It is uniformly generalized positive-lc: nonbad valuations already
have discrepancy at least $e$, the added divisor is effective Cartier
up to a rational multiple on this $\Q$-factorial model, and every bad
prime has received an increment at least $\sigma$.

Run its adjoint MMP on this Fano type model.  The Mori dream space
property and divisor MMP give a nef semiample transform of the big
$N$; see \cite[Corollary~1.3.1]{BCHM} and
\cite[Proposition~1.11]{HuKeel}.  Denote that transform by $N'$ on $Z$.
It is an effective nef big integral Weil divisor.  Negativity preserves
the uniform generalized discrepancy bound; forgetting the effective
boundary and b-nef part gives a uniform ordinary lc bound on $Z$.
The model remains of Fano type, so $-K_Z$ is big.  In particular
$N'-K_Z$ is big.  Theorem~1.1 of \cite{BirkarPolarised}, which permits
an integral Weil polarization, now makes $|qN'|$ birational for a
bounded positive integer $q$.  Negativity on a common resolution
transfers every such section to $qN$ on the original model: the
pullback of the source divisor dominates that of its transform.
This applies to the rational Cartier divisors represented by Weil
sections as well.  Finally $N\le k_0H$, so $|qk_0H|$ is birational and
$qk_0\le C_1/s$.

One can also obtain a bounded positive integer $u$ with $K_Z+uN'$ big.
This additional observation is useful when using the version of the
polarization result with both signs of $K_Z$.  Let $u_0\ge0$ be the
pseudo-effective threshold of $K_Z$ with respect to $N'$.  It is
rational by the Mori dream space property.  If $u_0=0$, bigness follows
for every positive $u$.  Otherwise take the adjoint MMP at $u_0$,
regarding $u_0\overline{N'}$ as b-nef big data.  Negativity retains a
uniform generalized lc bound, and the nef semiample threshold adjoint
is nonbig, hence defines a contraction of positive relative dimension.
On its geometric generic fibre, ordinary approximation with a small
ample reserve gives an effective uniformly positive-lc log Fano pair.
Theorem~1.1 of \cite{BirkarBAB} bounds the fibre.  For a bounded very
ample polarization, the degree of $-K$ is bounded above by adjunction
to a general section curve.  The transform of $N'$ has positive integral
degree there: it is effective and integral, and its bigness, preserved
by sections, excludes purely vertical support.  Relative adjoint
triviality gives $u_0\deg N'=\deg(-K)$, so $u_0$ is bounded.
Taking an integer $u>u_0$ proves the assertion.
\end{proof}

\subsection{Isolating the horizontal pivot}

Suppose the level is collapsing. Use
Lemma~\ref{b:scale-bad-primes} to choose an extraction $S'\to S$
containing the pivot $G_*$ as well as all bad primes, and put
$H_*=H(G_*)$. Set $b_*=1/4$, take $b\ll b_*$, and choose a fixed
rational $0<\nu\ll b_*$, sufficiently small also for the coefficient
bounds below. On $S'$ define
\[
 J=b_*G_*+\nu\sum_{\substack{P\ {\rm bad}\\P\ne G_*}}P.
\]
The coefficient at the pivot will determine the contraction threshold.
The smaller coefficients at the remaining bad primes retain a uniform
positive discrepancy bound without leaving another horizontal component.

\begin{lemma}\label{b:scale-collapse}
At a collapsing level there are a projective $\Q$-factorial Fano type
model $Y$ of $K$, a contraction $\pi:Y\to V$ with $\dim V<\dim S$,
an effective ample rational Cartier divisor $H_V$ on $V$, and a rational
$x\asymp s^{-1}$ with the following properties.  The pivot survives on
$Y$, and it is the only horizontal prime in the positive support of the
trace $H''$ of $H$.  Moreover
\[
 x=\frac{b_*}{H_*},\qquad H_*\asymp s,
 \qquad xH''-J''=x\pi^*H_V,
\]
with $J$ as just defined. On common higher models,
\begin{equation}\label{eq:B4}
                          H\ge\pi^*H_V.
\end{equation}
The datum with b-nef part
$\mathbf M_C=\mathbf M'+(x-d')\overline H$ and boundary obtained by
subtracting $J$ from the $A'$-boundary has effective uniformly
generalized positive-lc trace on $Y$, with total adjoint $x\pi^*H_V$.
The geometric generic fibres of $\pi$ form a bounded family.
\end{lemma}

\begin{proof}
\emph{The threshold and its lower bound.}
On the chosen extraction set
\[
 x=\min\{r:rH-J\text{ is pseudo-effective on }S'\}.
\]
Subtracting $J$ leaves the $A'$-boundary effective: its coefficient at
the pivot is at least $1-b-as$, and at the other indicated primes it
is greater than $1-e$.  The discrepancy becomes $A'+\overline J$.
Every bad prime has received a fixed positive increment, so this datum
is uniformly generalized $\delta'$-lc for a fixed $\delta'>0$.

The threshold $x$ is positive and finite, because $J$ is nonzero
effective and $H$ is big.  It is rational, and $xH-J$ is rationally
linearly equivalent to an effective divisor, by the Mori dream space
property.  Adding $J$ to such a representative and dividing by $x$
gives an effective representative of the pullback of $H$ on $S'$.
It descends as $H+\Div(f)/q$ to an effective representative on $S$,
by pushing its coefficients down.  Consequently weak essentiality at
the pivot gives
\begin{equation}\label{b:scale-pivot-lower}
                  \frac{b_*}{x}\le H_*+\eta s\le(1+\eta)s.
\end{equation}
It follows that $x\gtrsim s^{-1}$ and, after choosing $b$ small,
$x-d'\gtrsim x$.  Thus $\mathbf M_C$ is b-nef big.
Adding $(x-d')\overline H$ to the moduli changes the total adjoint to
$xH-J$ without changing $A'+\overline J$.

Run the adjoint MMP to a nef semiample transform on $Y$.  Its total
adjoint is the trace $xH''-J''$.  This divisor is not big, since $x$ is
the threshold.  Its semiample contraction $\pi:Y\to V$ therefore has
positive relative dimension.  The transformed boundary is effective,
and negativity preserves the uniform generalized $\delta'$-lc bound.
All the birational models used here remain $\Q$-factorial and of Fano
type.

\smallskip\noindent
\emph{Bounded fibres and the upper bound for the threshold.}
We next bound the degree of $xH''$ on a geometric generic fibre. A proper
normal generic fibre of a contraction in characteristic zero is
geometrically integral and normal.  Indeed its global functions are
$k(V)$, so $k(V)$ is algebraically closed in $k(Y)$; separability gives
the geometric assertions.  Canonical forms on the fibre are obtained
by dividing out a base volume form.  A resolution with snc data shows
that generic restriction and extension to an algebraic closure preserve
the required discrepancy bounds.  Apply
Lemma~\ref{b:ordinary-approximation} to the b-nef big part, reserving a
tiny ample class.  On the geometric generic fibre this gives an ordinary
effective $\delta'/2$-lc log Fano pair.  Theorem~1.1 of
\cite{BirkarBAB} therefore bounds the underlying fibre.  The theorem
applies over this algebraically closed characteristic-zero field; one
may also transport the data along an isomorphism of the algebraic
closure of $\C(V)$ with $\C$.

Fix a bounded very ample Cartier polarization on these fibres, and
write $\deg$ for intersection with its appropriate power.  If $L$
is this polarization in dimension $r\ge1$, a general
complete-intersection curve avoids the singular locus and gives
\[
             \deg(-K)\le (r-1)L^r+2.
\]
Thus this degree is bounded above.  The trace of any b-nef part has
nonnegative degree: upstairs add an arbitrarily small ample rational
class, take an effective representative, push its trace down, and then
let that class tend to zero.  On $Y$ the CY trace of $A^+$ still has
$B_0''\ge H''$ and is effective.  Every component of $J''$ is lc for
$A^+$, so its $B_0''$-coefficient is $1$ and $J''\le B_0''$.
The CY equality gives
\begin{equation}\label{b:scale-fibre-degree}
        x\deg H''=\deg J''\le\deg B_0''\le\deg(-K)\le C_3.
\end{equation}
Here traces are rational Cartier, since $Y$ is $\Q$-factorial.

The birational functions of Lemma~\ref{b:scale-birationality} remain
sections satisfying $\Div(f)+kH''\ge0$ after the non-extracting MMP.
If $sx$ were sufficiently large, \eqref{b:scale-fibre-degree} and
$k\le C_1/s$ would bound the degree of every possible polar divisor
on a geometric generic fibre by a number strictly less than $1$.
An integral nonzero effective divisor has positive integral degree,
so these polar divisors would vanish.  Properness would then put all
the functions in $k(V)$, contradicting their generation of $k(Y)$
and the positive relative dimension.  Thus $x\lesssim s^{-1}$.
Equation~\eqref{b:scale-pivot-lower} now also gives $H_*\gtrsim s$.

\smallskip\noindent
\emph{The pivot is the only positive horizontal component.}
Consider the coefficients of $H''-J''/x$.  At every positive coefficient
other than the pivot, this divisor has coefficient at least
$H''(P)/2\gtrsim s$.  Outside $J''$ this follows from the inductive
coefficient gaps and the extraction bounds.  At a bad prime use
$H(P)\ge s/2$, $1/x\le(1+\eta)s/b_*$, and choose $\nu$ sufficiently
small compared with $b_*$.  At the pivot its coefficient is at least
$-\eta s$, by \eqref{b:scale-pivot-lower}.  If horizontal, that prime
has bounded degree, since its coefficient in $B_0''$ is $1$.
The total degree of $H''-J''/x$ on a generic fibre is zero.
As $\eta\to0$, a horizontal nonpivot component would contribute at
least a fixed multiple of $s$, and the pivot could contribute only
$-O(\eta s)$.  This is impossible.  There must nevertheless be a
horizontal component in the positive support of $H''$, since otherwise
the birational functions would again be fibrewise regular.  Hence the
pivot survives and is the sole such horizontal component.  The degree
identity forces its coefficient in $H''-J''/x$ to be zero, proving
$x=b_*/H_*$.

\smallskip\noindent
\emph{The exact divisor on the base.}
It follows that $D_Y=xH''-J''$ is effective and vertical.  Semiampleness
gives $D_Y\sim_{\Q}\pi^*L$ for an ample rational Cartier divisor $L$
on $V$.  The rational function implementing this equivalence has no
horizontal poles, and hence lies in $k(V)$ by normality and properness
of the generic fibre.  Absorb its principal divisor into $L$ and divide
by $x$ to obtain an exact equality
$D_Y=x\pi^*H_V$.  Every strict base prime has a strict lift, so
$D_Y\ge0$ implies $H_V\ge0$.  Its rational linear class is ample.
Finally, on a common resolution, negativity gives the effective
rational Cartier b-divisor
\begin{equation}\label{b:scale-mmp-correction}
 \mathbf E_{\mathrm{MMP}}
       :=x\overline H-\overline J-x\overline{\pi^*H_V}\ge0.
\end{equation}
Since $J\ge0$, this proves \eqref{eq:B4}. The threshold datum after
the MMP has discrepancy
\[
 (A'+\overline J)+\mathbf E_{\mathrm{MMP}}
       =A'+x\overline H-x\overline{\pi^*H_V}.
\]
Thus its uniform positive lower bound holds on every prime over $Y$,
as required for the next step.
\end{proof}

\subsection{Passing the data to the base}

The contraction supplies $H_V$ and a uniformly positive discrepancy on
$Y$. We push this discrepancy to obtain $Q_{\rm next}$, and push $A^+$
separately to retain the lc places. Its bounded index will follow by
returning to the principal-moduli datum on the initial field.

\begin{proposition}\label{b:scale-inheritance}
At a collapsing level the construction of
Lemma~\ref{b:scale-collapse} supplies data at the next level satisfying
all the stated invariants.  More precisely, with the CY discrepancy
\[
                  C=A'+xH=a_K+\mathbf M_C,
\]
put
\begin{equation}\label{b:scale-next-data}
 \begin{aligned}
 A^+_{\rm next}&=\pi_\#A^+,&
 Q_{\rm next}&=\pi_\#(C-x\pi^*H_V),\\
 \lambda_{\rm next}&=x^{-1},&
 S_{\rm next}&=V,\qquad H_{\rm next}=H_V.
 \end{aligned}
\end{equation}
Then $\Lambda_{\rm next}=x^{-1}\pi_\#C$,
$\lambda_{\rm next}\asymp s\to0$, and $s_{\rm next}\gtrsim s$.
The CY moduli of $A^+_{\rm next}$ have a bounded b-Cartier multiple,
and $Q_{\rm next}$ has a uniform positive discrepancy bound on every
prime over $V$.
\end{proposition}

\begin{proof}
\emph{Discrepancy presentations and their lower bounds.}
By \eqref{b:scale-mmp-correction}, the threshold datum on $Y$ has
discrepancy exactly $C-x\pi^*H_V$. Therefore $C$ has effective horizontal trace and is
uniformly generalized klt over the generic point of $V$.  Its moduli
are b-nef big, so Proposition~\ref{b:push-presentation}, followed by
translation by the base divisor, gives a presentation for
$Q_{\rm next}$ with rational b-Cartier b-nef part and total adjoint
$xH_V$.

The CY datum $A^+$ has effective horizontal trace on $Y$.  Apply
Lemma~\ref{b:primary-minimization} to $A^+,C$.  For all sufficiently
small rational $t>0$, uniformly in the member and the base prime, a
lift computing the push of $(1-t)A^++tC$ also computes the push of
$A^+$.  On primary minimizers its value is
\[
       (1-t)A^+_{\rm next}(P)+t\frac{C(T)}{m_T}.
\]
A minimizer at any one such $t$ minimizes the second summand among
all primary minimizers.  Thus the push is exactly affine on a common
interval including $t=0$.  At positive $t$ the moduli are b-nef big
upstairs and the horizontal trace is klt, so
Proposition~\ref{b:push-presentation} supplies b-nef rational b-Cartier CY
moduli downstairs.  Two distinct positive rational parameters give
a common descent for their affine endpoint.  All smaller positive
parameters are nef on that descent; passing to the limit proves
b-nefness of the endpoint.  This constructs the required CY presentation
for $A^+_{\rm next}$.

Nonnegativity of $A^+_{\rm next}$ follows from \eqref{eq:B1}.
For effectivity, take a strict lift $F$ of a strict prime $P$ of $V$,
with restriction $mP$.  The effective trace upstairs bounds its
normalized discrepancy by $1/m\le1$, proving effectivity of the
pushed boundary.  For $A^+$ the stronger inequality
$B_0''\ge H''\ge\pi^*H_V$ gives
\[
 A^+_{\rm next}(P)\le\frac{A^+(F)}m
          \le\frac1m-H_V(P)\le1-H_V(P).
\]
Its boundary therefore dominates $H_V$.

To obtain a bound for $Q_{\rm next}$ on all valuations, apply
Lemma~\ref{b:ordinary-approximation} to the effective uniformly
generalized $\delta'$-lc threshold datum on the whole of $Y$.
Use an ordinary effective $\delta'/2$-lc representative with the same
relative rational log CY class and no larger discrepancies.
The Fano type property implies that $-K_Y$ is big over $V$.
The discriminant bound of \cite[Theorem~1.1]{BirkarFibrations}
therefore gives a uniform positive bound for the valuative push of
this ordinary pair, on every model of $V$.  Equation~\eqref{eq:B1}
then gives the same bound for $Q_{\rm next}$.  Notice that this
application uses the uniform lc bound on the entire total pair,
not just its generic fibre.

\smallskip\noindent
\emph{A model of the initial field and the uniform index.}
To bound the new index, we retain a model of the initial field throughout.  For each exceptional prime of $S'\to S$,
which is lc for $A^+$, attainment and composition in
Lemma~\ref{b:push-threshold} give a prime over $\mathcal Y$ that is
lc for $A_1^+$ and restricts to a positive multiple of the prescribed
prime.  Extract these primes, if they are not already strict, without
any others.  To justify this extraction, mix the effective ordinary
lc CY boundary of $A_1^+$ on $\mathcal Y$ with an effective klt log
Fano boundary, giving the latter a sufficiently small positive weight.
The mixture is log Fano and klt, and the requested discrepancies lie
strictly between $0$ and $1$.  Lemma~\ref{b:controlled-extraction}
applies, producing a $\Q$-factorial Fano type model $\mathcal Y'$.
The trace of $A_1^+$ remains effective, since each new prime has
$A_1^+$-discrepancy zero.  Every strict prime of $S'$ now has a strict
lift on $\mathcal Y'$; for the nonexceptional ones use the existing
morphism to $S$.  The same is true for every strict prime of $Y$,
because the MMP extracts none.

There need not yet be a morphism $\mathcal Y'\to Y$, so we construct
one without losing these properties.  Choose an ample Cartier divisor
$D_Y$ on $Y$, pull it back on a resolution of the rational map, and
let $D^*$ be its trace on $\mathcal Y'$.  It is rational Cartier.
The rational-function sections satisfy, in every nonnegative degree,
\begin{equation}\label{b:scale-auxiliary-sections}
 \{f\in K_1:\Div(f)+mD^*\ge0\}
 =\{f\in k(Y):\Div(f)+mD_Y\ge0\}.
\end{equation}
Indeed a function on the left has no poles on the normal proper generic
fibre of $\mathcal Y'\to S$: the divisor $D^*$ has coefficient zero
at every prime dominating $S$, since it comes from the field $k(S)$.
It is therefore a member of $k(S)=k(Y)$.  The strict lifts just
constructed test its order at every strict prime of $Y$, giving one
inclusion.  Conversely an effective Cartier divisor represented by a
section on $Y$ has effective pullback and trace on $\mathcal Y'$.
This also proves the assertion for the possibly nondivisible degrees,
by viewing the conditions as inequalities of rational divisors.

Run a divisor MMP for $D^*$ on the Fano type variety $\mathcal Y'$.
Its effective rational class is pseudo-effective, and its transform
is nef semiample.  Negativity preserves the divisible section systems
in \eqref{b:scale-auxiliary-sections}.  Their Proj is $Y$, so the
resulting model has a morphism to $Y$ inducing the given field
inclusion.  This morphism is a contraction: $k(Y)=k(S)$ is relatively
algebraically closed in $K_1$, and normality identifies its Stein
factor with $Y$.  The model remains $\Q$-factorial of Fano type, and
the trace of $A_1^+$ remains effective under non-extraction.
Composing with $Y\to V$ supplies the next auxiliary model.
Valuative composition shows that its push of $A_1^+$ is exactly
$A^+_{\rm next}$.

A klt log Fano boundary on this auxiliary model supplies an effective
generalized klt CY datum with b-nef big part: take minus its own
adjoint as the moduli divisor.  Apply
Corollary~\ref{b:principal-denominators} to this contraction and the
initial principal-moduli datum $A_1^+$.  Its dimension and denominator
$n$ are bounded independently of the level and member.  Thus a
bounded integer clears every discrepancy denominator of
$A^+_{\rm next}$.  On a smooth descent of its already constructed CY
moduli, the form discrepancies are integral, so the same integer
clears the coefficients of the moduli trace.  An integral divisor on
that smooth model is Cartier.  This proves the required bounded
b-Cartier multiple, without introducing any fibre-multiplicity factor.

\smallskip\noindent
\emph{Weak essentiality on the base.}
If $P$ is lc for $A^+_{\rm next}$,
choose a lift $T$, with restriction $m_TP$, computing the push of
$(1-t)A^++tC$ for a fixed sufficiently small rational $t>0$.
Lemma~\ref{b:primary-minimization} gives $A^+(T)=0$.  The formulas for
$C$, together with $x\asymp s^{-1}$, give discrepancy inequalities
\[
                       A^+\lesssim C,\qquad x\Lambda\lesssim C.
\]
For example, $C=(1-\gamma)A^++\gamma Q+(a+x)H$, all summands are
nonnegative, $1-\gamma$ is bounded below, and
$x\lambda/\gamma=xs/b$ is bounded when $Q$ is present.
Compare the mixture at $T$ with its value at a lift minimizing $C/m$.
Since $A^+(T)=0$, this proves
\begin{equation}\label{b:scale-essential-lift}
            \frac{\Lambda(T)}{m_T}
                  \le C_4\Lambda_{\rm next}(P).
\end{equation}
The factor $t^{-1}$ is harmless because $t$ was fixed uniformly.

If $H_V'\ge0$ is rationally linearly equivalent to $H_V$, apply its
rational principal change to $H$ through the field inclusion.
Inequality~\eqref{eq:B4} shows on a common resolution that the new
divisor is effective: it is the sum of $H-\pi^*H_V$ and $\pi^*H_V'$.
Pushing down gives an effective representative $H'\sim_{\Q}H$ on $S$.
Its change at $T$ is exactly
$m_T(P(H_V')-P(H_V))$.  Apply \eqref{b:scale-essentiality} at $T$
and then \eqref{b:scale-essential-lift}.  We obtain
\[
           P(H_V')\le P(H_V)+C_4\eta\Lambda_{\rm next}(P).
\]
Replacing $\eta$ by $C_4\eta$ preserves its convergence to zero.

\smallskip\noindent
\emph{Invariant-prime inequalities and strict coefficient gaps.}
For \eqref{eq:B2}, the first two inequalities follow immediately from
valuative minimization and \eqref{eq:B4}.  If an invariant prime has
restriction $jP$ at the old level and remains nontrivial at the next,
its full-multiplicity next $\Lambda$-value is at most
\[
 x^{-1}j\bigl((1-\gamma)A^+(P)+\gamma Q(P)+(a+x)H(P)\bigr).
\]
When $Q$ is present, \eqref{eq:B2} gives
\[
 j\gamma Q(P)\le\gamma+\frac bsC_0h_E.
\]
The constant $\gamma$ cancels the $-\gamma$ in the bound
$(1-\gamma)jA^+(P)\le(1-\gamma)(1-h_E)$.  Using $jH(P)\le h_E$ and
$x\asymp s^{-1}$, we obtain a bound
$x^{-1}+C'h_E$, exactly the next third inequality.  Initially the
same calculation simply omits the $Q$-term.

It remains to check the strict coefficient gaps and witnesses; they
cannot be inferred solely from the preceding inequalities.  For a
strict prime $P_V$ take a strict lift $F$ on $Y$ with restriction
$mP_V$.  The effective threshold boundary and the lower bound for
$Q_{\rm next}$ imply
\[
      0<\delta_{\rm next}\le Q_{\rm next}(P_V)
          \le\frac{(C-x\pi^*H_V)(F)}m\le\frac1m.
\]
Hence $m$ is uniformly bounded.  If $H_V(P_V)>0$ and
$A^+_{\rm next}(P_V)>0$, the lift is non-lc for $A^+$.  It is therefore
an old strict prime of $S$, lies outside $J$, and satisfies
$mH_V(P_V)=H''(F)\ge c_0$.  This proves the absolute positive gap.
If instead $P_V$ is lc and $H_V(P_V)>0$, the nonpivot coefficient bound
for $H''-J''/x$ gives $H_V(P_V)\gtrsim s$, since the pivot is horizontal
and $m$ is bounded.  Before any terminal replacement of the scale,
effectivity of the $Q_{\rm next}$-boundary gives
\[
 s_{\rm next}\le H_V(P_V)+\lambda_{\rm next}Q_{\rm next}(P_V)
                  \le H_V(P_V)+\lambda_{\rm next}.
\]
Since $\lambda_{\rm next}\asymp s$, this proves the required lower
bound by a fixed multiple of $s_{\rm next}$.

Finally suppose $H_V(P_V)=0$.  The same nonpivot coefficient bound
forces $H''(F)=0$.  All extracted primes have positive $H$-value, so
$F$ is an old strict prime.  Its old witness has $h_E=0$ and bounded
restriction multiplicity.  Restricting that witness to $k(V)$ only
multiplies this multiplicity by the bounded $m$, giving the next
witness.  The lower bound for $Q_{\rm next}$ also yields
$s_{\rm next}\gtrsim\lambda_{\rm next}\asymp s$ from the definition of
the scale.  All invariants have now been proved.
\end{proof}

\subsection{The terminal level}

At the last level the unmodified scale is bounded away from zero. We use
that lower bound to choose the constants below, while writing $s=1$ as
agreed above. No further field contraction is needed beyond the map to a
point.

\begin{lemma}\label{b:scale-terminal}
At the final level one may take $Y=S$ up to a small
$\Q$-factorialization, with target a point, and choose a fixed small
rational $x>0$ so that
\[
 C=A'+xH=a_K+\mathbf M_C,
 \qquad \mathbf M_C=\mathbf M'+(x-d')\overline H
\]
is effective and uniformly generalized klt.  Its b-nef big part has
a fixed positive $H$-reserve, its boundary dominates a fixed positive
multiple of $H$, and the underlying $Y$ belongs to a bounded family.
For every invariant prime with restriction $E|_{k(S)}=cP\ne0$,
\begin{equation}\label{b:scale-terminal-bound}
                         cC(P)\le1-\frac{h_E}{2}.
\end{equation}
\end{lemma}

\begin{proof}
Use $s=1$ and the actual lower bound for the unmodified scale to
choose $e$ as in Lemma~\ref{b:scale-bad-primes}.  There are no bad
primes, so $A'$, and hence $C\ge A'$, has a uniform positive discrepancy
bound.  Take $b\ll a$ and then a fixed small positive $x$.
Now $x-d'=x+a-b>0$ (initially it is $x+a$), giving a fixed positive
$H$-reserve.  The boundary of $C$ is
\[
           (1-\gamma)B_0+\gamma B_Q-(a+x)H
                  \ge(1-\gamma-a-x)H,
\]
with the usual initial convention.  Make the constants small enough
that the right side is a fixed positive multiple of $H$.

For the displayed estimate, the same constant-term cancellation as
in the inheritance proof gives
\[
 \begin{split}
 cC(P)&\le (1-\gamma)(1-h_E)
              +\gamma+bC_0h_E+(a+x)h_E\\
      &=1-(1-\gamma-bC_0-a-x)h_E.
 \end{split}
\]
Choose $\gamma+bC_0+a+x<1/2$; this is uniform since
$\gamma=b\lambda$ and $\lambda\to0$, or $\gamma=0$ initially.
This proves \eqref{b:scale-terminal-bound}.
Ordinary approximation with a tiny ample reserve gives an effective
ordinary uniformly positive-lc log Fano pair on $Y$.
Theorem~1.1 of \cite{BirkarBAB} bounds the underlying variety.
All assertions in dimension zero are interpreted vacuously.
\end{proof}

Starting with $S_1=S$, repeat Proposition~\ref{b:scale-inheritance}
whenever the scale tends to zero, and use
Lemma~\ref{b:scale-terminal} otherwise.  Each collapsing level strictly
lowers the dimension.  The construction therefore terminates after a
bounded number of levels, with contractions
\[
         Y_i\longrightarrow S_{i+1},\qquad
         k(Y_i)=k(S_i)=K_i,\qquad
         S_{f+1}=\Spec\C,
\]
where the $K_i$ form a nested tower of actual function fields.
Passing to a subsequence makes the number of levels and their dimensions
constant.  All positive constants fixed during this construction are
uniform along that subsequence.
